\section{基于Boosting的集成}
\label{sec:ensemble-boosting}

Bagging分别训练成员后才聚合；如果希望新成员针对当前预测误差学习，就需要先计算已有组合的输出。\term{Boosting}（提升）构造一个序贯加法模型：第$t$轮先观察已有函数$F_{t-1}$，再选择新成员及其贡献，以降低指定训练目标。AdaBoost、梯度提升及现代提升树共享这种依赖关系，但它们生成修正信号和近似目标的方式不同。

\subsection{前向逐步加法模型}
\label{sec:ensemble-stagewise-additive}

设$F_0$为初始预测，Boosting经过$T$轮得到
\begin{equation}
  F_T(\bm{x})
  =
  F_0(\bm{x})
  +
  \sum_{t=1}^{T}\alpha_t h_t(\bm{x}),
  \label{eq:ensemble-additive-model}
\end{equation}
其中$h_t$是第$t$个成员，$\alpha_t$控制其贡献。为使复杂度惩罚对应真正加入模型的函数，把本轮增量记为$f_t=\alpha_t h_t$。给定损失$\ell$，理想化的第$t$轮直接优化增量：
\begin{equation}
  f_t
  \in
  \argmin_{f\in\mathcal{F}}
  \sum_{i=1}^{n}
  \ell\!\left(
    y_i,
    F_{t-1}(\bm{x}_i)+f(\bm{x}_i)
  \right)
  +\Omega_t(f).
  \label{eq:ensemble-forward-stagewise}
\end{equation}
这里$\mathcal{F}=\{\alpha h:\alpha\in\mathbb{R},h\in\mathcal{H}\}$是由成员函数族$\mathcal{H}$生成的增量函数族，$\Omega_t$是相应复杂度惩罚，也可以取零。

树的叶数惩罚控制分区规模，叶输出范数惩罚控制本轮增量。范数应作用于$f$，使惩罚随实际加入组合的幅度变化；若仅惩罚$h$，它与$\alpha$的相互缩放会改变惩罚却保持预测。

逐步优化固定已有成员，减少了联合搜索变量，但最优范围仅限本轮增量。后文用梯度拟合、二阶目标和收缩构造近似更新。

不同方法将当前预测转换为不同的训练信号。AdaBoost改变训练样本的权重，使当前错分记录在下一轮具有更大影响；梯度提升把损失对当前预测的负梯度作为伪响应；二阶提升树还利用局部曲率近似叶权重和分裂收益。平方损失的负梯度恰为普通响应残差，其他损失则给出不同修正量。损失决定了下一轮成员拟合的对象。

\subsection{AdaBoost}
\label{sec:ensemble-adaboost}

\term{AdaBoost}（adaptive boosting）通过更新样本权重，使新分类器更多地处理当前组合反复错分的记录。考虑二分类标签$y_i\in\{-1,1\}$和成员$h_t(\bm{x})\in\{-1,1\}$。初始样本权重为$w_i^{(1)}=1/n$，第$t$轮成员的加权错误率为
\begin{equation}
  \epsilon_t
  =
  \sum_{i=1}^{n}
  w_i^{(t)}
  \mathbb{I}\!\left\{y_i\neq h_t(\bm{x}_i)\right\}.
  \label{eq:ensemble-adaboost-error}
\end{equation}
当$0<\epsilon_t<1/2$时，成员投票权重取
\begin{equation}
  \alpha_t
  =
  \frac{1}{2}
  \log\frac{1-\epsilon_t}{\epsilon_t},
  \qquad
  w_i^{(t+1)}
  =
  \frac{
    w_i^{(t)}
    \exp\!\left[-\alpha_t y_i h_t(\bm{x}_i)\right]
  }{Z_t},
  \label{eq:ensemble-adaboost-update}
\end{equation}
其中$Z_t$使新权重和为1。错分样本满足$y_i h_t(\bm{x}_i)=-1$，权重乘上$e^{\alpha_t}$；正确样本的权重乘上$e^{-\alpha_t}$。

若$\epsilon_t=0$，在这些正权重下当前成员已分对全部训练样本，便直接输出$h_t$并停止，避免将发散系数代入更新（参考\BookSectionRef{sec:boosting}{第二卷“基础、理论和可信性”的“二分类学习的可学习性”章“集成学习”节}）。其余情况下，最终分类器在预先固定平局规则后按$\operatorname{sign}(\sum_t\alpha_t h_t(\bm{x}))$决策。正的成员贡献要求当前加权错误率小于$1/2$；条件未满足时，可停止本轮、调整训练或采用相应变体。

图\ref{fig:ensemble-adaboost-weight-update}将分类错误与下一轮权重连接起来。当前分类结果先改变样本权重，再由归一化后的权重决定下一轮成员看到的训练分布；错分记录因而获得更大的后续影响。

\begin{figure*}[htbp]
  \centering
  \includegraphics[width=130.048mm]{figures/03-models/01-classic-models/generated/ensemble-adaboost-weight-update/figure.png}
  \caption{AdaBoost的样本权重更新。分类结果把样本分成正确与错分两路，两路缩放后的权重汇合并归一化，形成下一轮训练分布。实线表示常规训练与正确分类样本的更新路径，红色虚线强调错分样本被放大的路径。}
  \label{fig:ensemble-adaboost-weight-update}
\end{figure*}

这一更新也可以从指数损失解释。令$F(\bm{x})=\sum_t\alpha_t h_t(\bm{x})$，经验目标为
\[
  \sum_{i=1}^{n}\exp\!\left[-y_iF(\bm{x}_i)\right].
\]
固定$F_{t-1}$后选择$h_t$与$\alpha_t$，在$0<\epsilon_t<1/2$时可以导出式\eqref{eq:ensemble-adaboost-update}。样本重加权由指数损失的逐步优化导出，因此其作用可以通过该目标分析\scite{freund1997boosting,friedman2000adaboost}。

持续错分的记录获得更大权重，后续成员因而可能恢复稀少的边界结构，也可能追逐错误标签。若成员在每轮带权分布下都优于随机猜测，可以控制训练错误的下降；新样本风险还与投票间隔和函数类条件有关（参考\BookSectionRef{sec:boosting}{第二卷“基础、理论和可信性”的“二分类学习的可学习性”章“集成学习”节}）。实际训练需同时检查错误记录的来源、训练目标和独立测试表现。

多分类AdaBoost需要重新定义成员错误与类别得分，回归变体则要用响应误差决定样本权重和预测聚合。这些变体沿用序贯重加权，但成员错误、系数与聚合需要按各自任务重新推导。评价分别采用分类或回归的输出和损失。

\subsection{梯度提升机与GBDT}
\label{sec:ensemble-gradient-boosting}

改变任务损失后，也需要改变修正信号。对于关于预测可微的损失，其负梯度指出当前预测降低经验目标的方向。\term{梯度提升机}（gradient boosting machine, GBM）用基学习器逼近这个负梯度方向，从而把序贯加法扩展到更一般的损失\scite{friedman2001gradientboost}。

通常从使训练损失最小的常数预测开始：
\begin{equation}
  F_0(\bm{x})=c_0,
  \qquad
  c_0\in\argmin_c\sum_{i=1}^{n}\ell(y_i,c).
  \label{eq:ensemble-gbm-initialization}
\end{equation}
平方损失下$c_0$是响应均值。这里要求存在有限的最小化点；若某种损失与数据组合使最优常数趋于无穷，就需约束预测范围或另定有限初始化。

在第$t$轮，对每个训练样本计算\term{伪残差}（pseudo-residual），即损失在当前预测处的负导数：
\begin{equation}
  r_i^{(t)}
  =
  -
  \left.
  \frac{\partial \ell(y_i,z)}{\partial z}
  \right|_{z=F_{t-1}(\bm{x}_i)}.
  \label{eq:ensemble-pseudo-residual}
\end{equation}
随后用平方误差把这些修正量拟合到成员函数族：
\begin{equation}
  h_t\in\argmin_{h\in\mathcal{H}}
  \sum_{i=1}^{n}\bigl(r_i^{(t)}-h(\bm{x}_i)\bigr)^2.
  \label{eq:ensemble-gbm-gradient-fit}
\end{equation}
此处平方误差衡量成员对负梯度的逼近，即使原任务采用其他损失也使用这一拟合对象。贪心树生长通常给出子问题的近似解。

固定$h_t$后，可以按原损失做整体线搜索，在最小值存在时取
\begin{equation}
  \alpha_t\in\argmin_\alpha
  \sum_{i=1}^{n}
  \ell\!\left(y_i,F_{t-1}(\bm{x}_i)+\alpha h_t(\bm{x}_i)\right).
  \label{eq:ensemble-gbm-line-search}
\end{equation}
此处$\alpha_t$是整棵树或整个成员共享的一个标量。加入预先选择的收缩系数$\eta\in(0,1]$后，
\begin{equation}
  F_t(\bm{x})
  =
  F_{t-1}(\bm{x})
  +
  \eta\alpha_t h_t(\bm{x}).
  \label{eq:ensemble-gradient-boosting-update}
\end{equation}


当$\ell(y,z)=\tfrac12(y-z)^2$时，式\eqref{eq:ensemble-pseudo-residual}给出$y_i-F_{t-1}(\bm{x}_i)$，恰好是普通残差。分类时则要先明确得分与概率的关系。例如二分类采用$y_i\in\{0,1\}$，令$F$表示对数几率（logit），相应概率与二项对数损失为
\[
  p=\sigma(F)=\frac{1}{1+e^{-F}},\qquad
  \ell(y,F)=\log(1+e^F)-yF.
\]
由于$\partial\ell/\partial F=\sigma(F)-y$，第$t$轮的伪残差为
\[
  r_i^{(t)}=y_i-p_i^{(t-1)},\qquad
  p_i^{(t-1)}=\sigma\!\left(F_{t-1}(\bm{x}_i)\right).
\]


例如$y_i=1$而当前概率为$0.3$时，伪残差为$0.7$，表示损失对得分的修正信号。成员拟合后将增量加到$F$，再经Sigmoid换算概率；修正量与概率因此处于不同数值空间。

若$h_t$采用回归树，就得到\term{梯度提升决策树}（gradient boosting decision tree, GBDT）。除了整体线搜索，还可以保留平方拟合得到的叶区域$\bm{R}_1^{(t)},\ldots,\bm{R}_{L_t}^{(t)}$，再按原损失分别求每片叶的增量。当各叶内存在有限最小化点时，
\begin{equation}
  \begin{aligned}
    v_j^{(t)}
      &\in\argmin_v
        \sum_{i:\,\bm{x}_i\in\bm{R}_j^{(t)}}
        \ell\!\left(y_i,F_{t-1}(\bm{x}_i)+v\right),\\
    F_t(\bm{x})
      &=F_{t-1}(\bm{x})
        +\eta\sum_{j=1}^{L_t}
        v_j^{(t)}\mathbb{I}\{\bm{x}\in\bm{R}_j^{(t)}\}.
  \end{aligned}
  \label{eq:ensemble-gbm-leaf-update}
\end{equation}
$L_t$表示叶数，各叶分别选择$v_j^{(t)}$并替换原树输出；这与整棵树共用一个$\alpha_t$是两种更新方案。有限叶内最小值不存在时，需要限制输出或增加正则后再更新。回归树在这里承担连续修正量的拟合，最终任务也可以是分类。

第\ref{sec:regression-quantiles}节的分位数损失$\ell(y,z)=\rho_\tau(y-z)$在$z=y$处不可微，修正信号需改为次梯度。分位数提升可改用负次梯度，例如约定
\begin{equation}
  r_i^{(t)}
  =\tau-\mathbb{I}\{y_i<F_{t-1}(\bm{x}_i)\},
  \qquad 0<\tau<1.
  \label{eq:ensemble-quantile-pseudo-residual}
\end{equation}
等号处的负次梯度可取$[\tau-1,\tau]$内任意值，这里选端点$\tau$。常数初始化取样本$\tau$分位数，仍用平方误差拟合上述修正量；逐叶更新则最小化原分位数损失，故$v_j^{(t)}$取该叶内普通残差$y_i-F_{t-1}(\bm{x}_i)$的$\tau$分位数，而不是伪残差的均值。

学习率、轮数与单树复杂度共同决定加法函数。较小学习率缩小每轮变化，通常需要更多轮才能达到相近训练目标；较深的树可以在单轮表达更高阶交互，也会扩大每轮容量。行采样和列采样进一步引入随机性，早停则依据独立验证轨迹选择轮数。选择学习率时需同时调整轮数，并在给定树复杂度、采样和早停协议下比较。

下一轮使用当前组合的梯度，GBM各轮因而按序执行。单轮内部仍可并行统计不同特征或候选分裂。训练损失描述经验拟合，停止轨迹描述求解，测试风险描述未来预测；若轮数、损失、树结构和早停规则都根据同一验证集反复调整，还要把这种选择纳入最终评价。

\subsection{XGBoost}
\label{sec:ensemble-xgboost}

经典梯度提升用一阶负梯度给出拟合方向。\term{XGBoost}在正则化加法树目标中进一步使用二阶局部近似，并围绕稀疏输入、近似分裂和内存访问设计训练系统\scite{chen2016xgboost}。二阶目标给出候选树的评分，分裂算法搜索结构，系统设计组织统计量的计算与访问。三层共同形成实现，但回答不同问题。

设第$t$棵树为$f_t$，损失在当前预测$F_{t-1}(\bm{x}_i)$处的一阶导与二阶导分别为$g_i^{(t)}$和$s_i^{(t)}$。去掉与$f_t$无关的常数后，本轮目标近似为
\begin{equation}
  \widetilde{\mathcal{L}}^{(t)}
  =
  \sum_{i=1}^{n}
  \left[
    g_i^{(t)}f_t(\bm{x}_i)
    +\frac{1}{2}s_i^{(t)}f_t^2(\bm{x}_i)
  \right]
  +\gamma L
  +\frac{\lambda}{2}\sum_{j=1}^{L}v_j^2,
  \label{eq:ensemble-xgboost-objective}
\end{equation}
其中$L$是叶数，$v_j$是第$j$片叶的输出，$\gamma$惩罚新增叶，$\lambda$约束叶权重。若$I_j$表示落入第$j$片叶的样本索引，记
\[
  G_j=\sum_{i\in I_j}g_i^{(t)},
  \qquad
  S_j=\sum_{i\in I_j}s_i^{(t)}.
\]
在$S_j+\lambda>0$时，固定树结构的最优叶权重为
\begin{equation}
  v_j^\star=-\frac{G_j}{S_j+\lambda}.
  \label{eq:ensemble-xgboost-leaf-weight}
\end{equation}
把一个父叶分成左右子叶时，相应近似目标下降量为
\begin{equation}
  \operatorname{Gain}
  =
  \frac{1}{2}
  \left[
    \frac{G_{\mathrm{L}}^2}{S_{\mathrm{L}}+\lambda}
    +
    \frac{G_{\mathrm{R}}^2}{S_{\mathrm{R}}+\lambda}
    -
    \frac{(G_{\mathrm{L}}+G_{\mathrm{R}})^2}
         {S_{\mathrm{L}}+S_{\mathrm{R}}+\lambda}
  \right]
  -\gamma.
  \label{eq:ensemble-xgboost-split-gain}
\end{equation}
式\eqref{eq:ensemble-xgboost-split-gain}把当前损失的梯度、曲率与树复杂度放在同一个分裂准则中。增益评价局部二阶目标在当前分裂下的变化；贪心生长只搜索逐步形成的结构。

图\ref{fig:ensemble-gradient-xgboost}把一阶梯度提升与XGBoost的二阶树评分放在同一计算链中。左图用负梯度定义拟合对象，右图从候选区域的梯度统计计算叶权重和分裂增益；二阶统计扩展了本轮目标的近似信息，但没有改变逐轮加法更新的外层结构。

\begin{figure*}[htbp]
  \centering
  \includegraphics[width=152.739mm]{figures/03-models/01-classic-models/generated/ensemble-gradient-xgboost/figure.png}
  \caption{梯度提升与XGBoost中的本轮修正。左图中，损失对当前预测的负梯度形成伪残差，新成员拟合该信号后更新组合；右图中，XGBoost按候选分裂聚合一阶导与二阶导，再由同一组叶统计计算叶权重和分裂增益。图中单叶权重$w^\ast$对应正文的$v_j^\star$，$G$与$H$分别表示叶内一阶导和二阶导之和。复杂度项$\gamma T+\lambda\lVert\bm{w}\rVert_2^2/2$中的$T$为叶数，$\bm{w}$汇集各叶输出；$\lambda$限制叶权重，$\gamma$惩罚新增叶。该图展示当前候选树的局部评分。}
  \label{fig:ensemble-gradient-xgboost}
\end{figure*}

精确分裂搜索遍历排序后的候选特征值；近似搜索则先用带权分位数草图构造候选切分，再汇总梯度统计。稀疏感知搜索为缺失值学习默认方向，列采样减少每轮需要检查的特征。块式存储、缓存感知访问、并行统计和外存计算改变的是实现代价，不改变式\eqref{eq:ensemble-xgboost-objective}所定义的学习对象。这些设计在特定硬件和数据规模上的时间收益可通过完整训练测量确认。

\subsection{LightGBM}
\label{sec:ensemble-lightgbm}

逐值分裂搜索需要遍历候选取值，样本和特征增多会扩大这一工作量。\term{LightGBM}以直方图训练为基础，把连续特征映射到有限个桶，并在桶级别累计梯度和曲率。构造直方图仍要访问节点样本，但候选分裂只需扫描桶；已知父节点和一个子节点的直方图时，还可以用差分得到另一个子节点的统计量。分桶减少存储和候选数，也引入离散化误差。

树的生长顺序进一步改变计算分配。按层生长同时扩展某一深度上的叶，按叶生长则每次选择当前增益最大的叶继续分裂。在相同叶数下，后者把容量集中到训练目标下降较大的局部区域，可能形成深度不均衡的树。叶数、最小叶样本量和深度共同限制这种局部扩展，训练收益再由验证或测试检查其样本外效果。

LightGBM论文进一步提出两种针对大规模数据的近似。\term{单边梯度采样}（gradient-based one-side sampling, GOSS）保留绝对梯度较大的样本，从绝对梯度较小的样本中随机抽取一部分，并在分裂统计中重新加权后者。例如，在小梯度组内等概率抽取一半样本，就把被抽中者对梯度和的贡献乘二，以补偿未采样部分。一般地，样本以概率$p>0$被保留时，其贡献乘以$1/p$；给定本轮梯度和候选分区，这种逆包含概率加权可无偏估计相应的梯度和，平方、比值及择优得到的增益则是非线性量，其偏差需另行分析。

\term{互斥特征捆绑}（exclusive feature bundling, EFB）把很少同时非零的稀疏特征编码到较少的组合特征中，近似减少有效特征数。GOSS、EFB和直方图分别减少样本统计、特征表示和阈值候选，可按各自近似对象评价误差与计算节省\scite{ke2017lightgbm}。

图\ref{fig:ensemble-lightgbm-approximations}按近似对象并列这三种机制。直方图压缩候选阈值，GOSS压缩参与分裂统计的样本，EFB压缩稀疏特征表示；这三个对象分别对应数据访问、有效特征数与候选扫描的成本。

% Flush preceding margin references before the full-width illustration.
\clearpage
\begin{figure*}[htbp]
  \centering
  \includegraphics[width=130.048mm]{figures/03-models/01-classic-models/generated/ensemble-lightgbm-approximations/figure.png}
  \caption{LightGBM中三种作用位置不同的计算近似。直方图把连续取值映射到有限桶，GOSS保留$|g_i|$较大的样本，并抽取、重加权部分$|g_i|$较小的样本，EFB把很少同时非零的特征编码为较少的捆绑特征。图中点数、桶数与非零位置均为机制示意。}
  \label{fig:ensemble-lightgbm-approximations}
\end{figure*}

这些机制的效果依赖梯度分布、特征稀疏性、桶数、冲突处理和硬件实现。比较速度和任务损失时，需要记录数据规模、硬件、桶数与采样设置，并固定搜索预算及早停规则。论文实验的结论据此限定于相应条件。

\subsection{CatBoost}
\label{sec:ensemble-catboost}

高基数类别特征若全部展开为独热变量，会增加特征规模；把类别直接映射为整数，又会引入原本不存在的次序。另一种做法是用同类别样本的响应均值编码，但若样本自己的$y_i$参与其编码，模型就能从输入中间接看到待预测标签。训练输入含有待预测响应、部署输入却没有，这种构造差异称为\term{目标泄漏}（target leakage）。

\term{CatBoost}用随机排列建立\term{有序目标统计}（ordered target statistic）。设$c_i$为样本$i$的某个类别值，用$\pi(j)$表示样本$j$在随机排列中的位置；对当前样本只使用排列中位于它之前且类别相同的记录：
\begin{equation}
  z_i
  =
  \frac{
    \sum_{j:\,\pi(j)<\pi(i)}
    \mathbb{I}\{c_j=c_i\}y_j
    +ap
  }{
    \sum_{j:\,\pi(j)<\pi(i)}
    \mathbb{I}\{c_j=c_i\}
    +a
  }.
  \label{eq:ensemble-catboost-ordered-statistic}
\end{equation}
其中$p$是先验统计量，$a>0$控制平滑。前缀统计排除了样本自身的标签，部署时可从已观测的训练记录形成同类别统计。低频类别仍会产生较大估计波动，平滑、多个排列和类别组合需要共同考虑。

图\ref{fig:ensemble-catboost-ordered-statistics}固定当前样本的排列位置，检查哪些标签进入编码。计算当前类别$A$样本的编码时，只允许排列前缀中同为类别$A$的标签进入统计量；其他类别、当前标签和排列后缀都不参与这次编码。

\begin{figure*}[htbp]
  \centering
  \includegraphics[width=150.199mm]{figures/03-models/01-classic-models/generated/ensemble-catboost-ordered-statistics/figure.png}
  \caption{CatBoost有序目标统计的前缀约束。这里$\pi(j)$表示样本$j$在随机排列中的位置；计算当前样本$i$的编码时，只有满足$\pi(j)<\pi(i)$且$c_j=c_i$的标签通过实线汇入统计量。其他类别、当前标签和后续记录均不连接到统计量，表示它们在本次编码中被隔离。样本索引与类别符号仅用于说明机制。}
  \label{fig:ensemble-catboost-ordered-statistics}
\end{figure*}

有序提升把相同的前缀思想用于梯度估计。计算当前样本的预测与梯度时，采用由其前缀记录拟合的相应模型，使该样本的标签与当前预测拟合分离，以减少这一重复使用带来的预测偏移。实际算法通过多个排列和前缀模型实现这一约束，训练状态与计算量也比普通单序列提升更复杂\scite{prokhorenkova2018catboost}。

CatBoost常使用\term{对称树}（symmetric tree），即同一深度的所有节点采用相同分裂条件。深度为$d_{\mathrm{tree}}$的完整对称树至多有$2^{d_{\mathrm{tree}}}$片叶，路径可以编码为固定长度的二进制选择。这种结构便于批量预测，也限制了单棵树可表达的分区形状。类别编码改变输入，有序提升改变梯度估计所用模型，对称树限制分区结构；三者的收益与代价应按各自对象评价。

\subsection{模型对比与总结}
\label{sec:ensemble-boosting-comparison}

表\ref{tab:ensemble-boosting-comparison}先固定共同的序贯加法结构，再比较各方法修改哪一层。AdaBoost由加权分类错误生成样本权重；GBM由损失梯度生成伪残差；XGBoost对正则化树目标作二阶近似；LightGBM进一步压缩分裂统计的计算；CatBoost重点处理类别统计与序贯估计中的目标泄漏。具体系统可同时使用多种训练选项，表中按代表机制比较，实际配置仍需单独记录。

\begin{table*}[htbp]
  \centering
  \small
  \begin{tabularx}{\linewidth}{@{}>{\raggedright\arraybackslash}p{19mm}*{6}{>{\raggedright\arraybackslash}X}@{}}
    \toprule
    方法 & 本轮修正信号 & 典型成员 & 目标近似 & 树生长或分裂 & 类别特征重点 & 主要约束 \\
    \midrule
    AdaBoost & 加权错分 & 树桩或浅分类树 & 指数损失下的解析更新 & 由成员算法决定 & 需另行编码 & 对持续错分记录敏感 \\
    GBM／GBDT & 负梯度伪残差 & 回归树 & 一阶函数方向与叶内优化 & 通常贪心树生长 & 需随实现处理 & 学习率、轮数与树复杂度耦合 \\
    XGBoost & 一阶导与二阶导 & 正则化回归树 & 二阶局部目标 & 精确或分位数候选；稀疏默认方向 & 依实现与版本而定 & 轮次串行，近似与系统配置影响成本 \\
    LightGBM & 桶级梯度统计 & 直方图回归树 & 通常使用二阶统计 & 按叶生长；直方图候选 & 支持类别分组，机制依实现 & 分桶、GOSS与EFB引入近似条件 \\
    CatBoost & 有序统计与有序梯度 & 对称树 & 梯度或二阶统计 & 同层共享分裂条件 & 有序目标统计与组合 & 多排列和前缀状态增加训练成本 \\
    \bottomrule
  \end{tabularx}
  \caption{Boosting代表方法的机制比较。GBM是一般框架，GBDT指使用回归树成员的实例；XGBoost、LightGBM和CatBoost是树提升的具体模型与系统，均属于Boosting框架的具体构造。}
  \label{tab:ensemble-boosting-comparison}
\end{table*}

% 跨栏表之后的扩展段需与三篇完整出处共同容纳，不能只按正文行数分页。
\Needspace{23\baselineskip}
随机梯度提升在每轮使用训练子样本，DART随机暂时移除部分已有树后计算修正，LogitBoost则从加法Logistic回归推导更新\scite{friedman2002stochastic,rashmi2015dart,friedman2000adaboost}。三种扩展分别改变抽样、当前组合及损失近似。Boosting能够沿训练目标定向增加容量，却仍采用预先规定的加法组合形式；若已有成员来自不同模型族，并希望从它们的样本外表现中学习组合方式，就进入Stacking的范围。
